\begin{table}[H]
\centering
\caption{Matriz de impacto y efectos de largo plazo estimados}
\label{tab:impacto}
\small
\begin{tabular}{@{}lrrrrr@{}}
\toprule
Variable & $\varepsilon^{ISE}$ & $\varepsilon^{\pi}$ & $\varepsilon^{\pi^e}$ & $\varepsilon^{TPM}$ & $\varepsilon^{ITCR}$ \\
\midrule
\multicolumn{6}{@{}l}{\textit{A. Impacto: $\hat H$}} \\
ISE & $1.032$ & $0$ & $0$ & $0$ & $0$ \\
Inflación & $-0.024$ & $0.171$ & $0$ & $0$ & $0$ \\
Expectativas & $-0.024$ & $0.045$ & $0.176$ & $0$ & $0$ \\
TPM & $-0.004$ & $0.053$ & $0.002$ & $0.161$ & $-0.017$ \\
ITCR & $-0.187$ & $-0.129$ & $0.035$ & $-0.040$ & $2.289$ \\
\addlinespace
\multicolumn{6}{@{}l}{\textit{B. Largo plazo: $(I - \hat A(1))^{-1}\hat H$}} \\
ISE & $0.383$ & $-0.470$ & $-0.130$ & $-0.390$ & $0.021$ \\
Inflación & $0.059$ & $0.553$ & $0.075$ & $-0.116$ & $0.326$ \\
Expectativas & $0.126$ & $1.734$ & $1.438$ & $-0.171$ & $0.410$ \\
TPM & $1.140$ & $8.974$ & $3.371$ & $2.816$ & $1.256$ \\
ITCR & $-0.122$ & $-0.827$ & $-1.165$ & $0$ & $2.381$ \\
\bottomrule
\end{tabular}
\notascuadro{\textit{Nota.} Cada columna es un choque estructural de una desviación estándar y cada fila, la respuesta de una variable. Panel A: respuesta en el mes del choque; los ceros son las restricciones de corto plazo. Panel B: efecto acumulado de largo plazo; el cero en la fila del ITCR y la columna de la TPM es la restricción de largo plazo. El choque monetario (columna TPM) se obtiene rotando las dos últimas columnas de la descomposición de Cholesky un ángulo de 5.9 grados. Un choque de una desviación estándar eleva la TPM en 0.161 pp; las respuestas al impulso se reescalan a un alza de 25 pb.}
\end{table}
